\chapter*{Phụ lục}
\addcontentsline{toc}{chapter}{Phụ lục}

\section*{A. Cấu trúc thư mục dự án}

\begin{verbatim}
thesis-smart-home-gas-detection/
├── application/
│   ├── backend/src/
│   │   ├── controllers/dashboard.controller.ts
│   │   ├── services/influx.service.ts
│   │   ├── services/postgres.service.ts
│   │   └── services/kafka.consumer.ts
│   └── frontend/views/dashboard.ejs
├── communication/bridges/mqtt_kafka_bridge.py
├── device/simulator/sensor_simulator.py
├── infrastructure/
│   ├── grafana/provisioning/
│   ├── influxdb/
│   ├── mosquitto/mosquitto.conf
│   └── postgresql/init/01_schema.sql
├── processing/
│   ├── ml/
│   │   ├── inference/forecaster_inference.py
│   │   ├── inference/rl_inference.py
│   │   ├── lstm/train_forecaster.py
│   │   └── rl/train_rl.py, gas_env.py
│   ├── scripts/benchmark.py
│   └── spark-streaming/jobs/stream_processor.py
├── docker-compose.yml
└── README.md
\end{verbatim}

\section*{B. Hướng dẫn cài đặt và chạy hệ thống}

\subsection*{B.1. Yêu cầu hệ thống}
\begin{itemize}
  \item Docker Engine 24+ và Docker Compose v2
  \item Python 3.11+ (nếu huấn luyện mô hình ngoài Docker)
  \item RAM tối thiểu 8 GB (khuyến nghị 16 GB cho Spark)
  \item Ổ đĩa trống tối thiểu 10 GB
\end{itemize}

\subsection*{B.2. Huấn luyện mô hình (tùy chọn)}

Nếu muốn huấn luyện mô hình trước khi chạy hệ thống:

\begin{verbatim}
pip install -r processing/requirements-processing.txt

# Tái tạo thí nghiệm LSTM mô phỏng trong Chương 4
python processing/ml/lstm/train_forecaster.py \
    --hours 8 --epochs 20 --seed 42

# Tái tạo thí nghiệm PPO trong Chương 4
python processing/ml/rl/train_rl.py --steps 400000 --seed 42 --n-envs 4
\end{verbatim}

Nếu bỏ qua bước này, hệ thống sẽ tự động sử dụng fallback mechanism.

\subsection*{B.3. Chạy benchmark đánh giá}

\begin{verbatim}
python processing/scripts/benchmark.py --hours 4 --seed 42 --seeds 3
python processing/scripts/benchmark_scenarios.py --seeds 3 --seed-start 42
\end{verbatim}

\subsection*{B.4. Khởi động toàn bộ hệ thống}

\begin{verbatim}
cp .env.example .env          # Sao chép cấu hình mẫu
# Chỉnh sửa .env nếu cần (Telegram token, credentials...)
docker compose up --build     # Build và khởi động tất cả dịch vụ
\end{verbatim}

Sau khi khởi động:
\begin{itemize}
  \item Bảng điều khiển web: \texttt{http://localhost:8080}
  \item Grafana: \texttt{http://localhost:3001} (admin / admin123456)
  \item InfluxDB UI: \texttt{http://localhost:8086}
\end{itemize}

\section*{C. Cấu hình biến môi trường}

\begin{table}[H]
\centering
\caption{Các biến môi trường quan trọng}
\begin{tabular}{lll}
\toprule
\textbf{Biến} & \textbf{Mặc định} & \textbf{Mô tả} \\
\midrule
\texttt{MQTT\_BROKER\_HOST}    & mosquitto & Host của MQTT broker \\
\texttt{MQTT\_BROKER\_PORT}    & 1883      & Port MQTT \\
\texttt{KAFKA\_BOOTSTRAP\_SERVERS} & kafka:9092 & Kafka brokers \\
\texttt{INFLUXDB\_URL}         & http://influxdb:8086 & URL InfluxDB \\
\texttt{INFLUXDB\_TOKEN}       & --        & Access token InfluxDB \\
\texttt{POSTGRES\_URL}         & --        & Connection string PostgreSQL \\
\texttt{TELEGRAM\_BOT\_TOKEN}  & --        & Bot token từ BotFather \\
\texttt{TELEGRAM\_CHAT\_ID}    & --        & Chat ID nhận cảnh báo \\
\texttt{SIM\_LEAK\_PROB\_PER\_MIN} & 0.05  & Xác suất rò rỉ mỗi phút \\
\texttt{SIM\_PUBLISH\_HZ}      & 1.0       & Tần suất gửi dữ liệu cảm biến \\
\bottomrule
\end{tabular}
\end{table}
